\documentclass[12pt,a4paper]{ctexart}

\usepackage[top=2.5cm, bottom=2.5cm, left=2.8cm, right=2.8cm]{geometry}
\usepackage{amsmath,amssymb,amsthm}
\usepackage{booktabs}
\usepackage{enumitem}
\usepackage[colorlinks=true,linkcolor=blue,citecolor=blue]{hyperref}
\usepackage{xcolor}
\usepackage{algorithm}
\usepackage{algpseudocode}

\newcommand{\tdet}{t_{\text{det}}}
\newcommand{\Teff}{T_{\text{eff}}}
\newcommand{\GP}{\mathcal{GP}}
\newcommand{\EI}{\mathrm{EI}}
\newcommand{\R}{\mathbb{R}}
\newcommand{\N}{\mathcal{N}}
\newcommand{\X}{\mathbf{X}}
\newcommand{\y}{\mathbf{y}}
\newcommand{\K}{\mathbf{K}}
\newcommand{\kstar}{\mathbf{k}_*}
\newcommand{\LS}{\mathbf{L}}
\newcommand{\alphavec}{\boldsymbol{\alpha}}

\title{\textbf{问题二：锚点搜索 + 贝叶斯优化求解最优单弹部署}\\[5pt]
       \large 网格搜索 $\to$ 坐标下降 $\to$ GP + EI —— 原理推导、算法实现与结果分析}
\author{}
\date{\today}

\begin{document}

\maketitle

\section{问题描述与数学建模}

\subsection{场景与决策变量}

无人机 FY1 从 $\mathbf{P}_{\text{FY1}} = (17800, 0, 1800)$ 出发，以恒定速度 $v \in [70, 140]\ \mathrm{m/s}$ 在高度 $H = 1800\ \mathrm{m}$ 水平直线飞行。收到任务后，FY1 确定一个不变的飞行方向 $\theta$（航向角，以 $+X$ 轴为 $0^\circ$ 逆时针度量）和巡航速度 $v$，此后的飞行中二者均不再改变。投放 1 枚烟幕干扰弹，干扰导弹 M1。

因为无人机在整个飞行阶段航向和速度不变，问题一中由投放时刻 $t_{\text{rel}}$ 和延时 $t_{\text{delay}}$ 描述的正向参数化等价于直接指定\textbf{起爆时刻 $\tdet$}和\textbf{起爆点 $\mathbf{A} = (A_x, A_y, A_z)$}。无人机飞行速度由起爆点与起点的水平距离反推，投放时刻由起爆时刻减去自由落体时间得到。因此 4 维决策向量为：

\begin{equation}
\boxed{\mathbf{x} = [\,\tdet,\; A_x,\; A_y,\; A_z\,] \in \R^4}
\end{equation}

$\tdet$ 是烟幕弹从无人机起爆的时刻（单位：秒），$(A_x, A_y, A_z)$ 是起爆瞬间烟幕球心的三维空间坐标（单位：米）。给定 $\mathbf{x}$ 后，无人机的飞行参数由以下逆向公式唯一确定：

\begin{align}
v(\mathbf{x}) &= \frac{\sqrt{(A_x - 17800)^2 + (A_y - 0)^2}}{\tdet} \label{eq:v} \\
t_{\text{fall}}(\mathbf{x}) &= \sqrt{\frac{2(1800 - A_z)}{g}}, \quad g = 9.8\ \mathrm{m/s^2} \label{eq:tfall} \\
t_{\text{rel}}(\mathbf{x}) &= \tdet - t_{\text{fall}} \label{eq:trel}
\end{align}

式(\ref{eq:v})的含义：无人机从起点 $\mathbf{P}_{\text{FY1}}$ 以恒定速度 $v$ 沿直线飞行 $\tdet$ 秒后，其水平投影恰好到达起爆点 $(A_x, A_y)$ 的正上方。式(\ref{eq:tfall})是自由落体时间——炸弹从 1800m 高度释放后下落到 $A_z$ 所需的时间。式(\ref{eq:trel})要求投放时刻非负，即 $t_{\text{rel}} \geq 0$。

\subsection{约束条件的物理含义}

\begin{enumerate}[label=(C\arabic*), leftmargin=*]
    \item \textbf{速度限制} $v \in [70, 140]$：FY1 的飞行速度受动力系统限制，低于 70m/s 将失速，高于 140m/s 超出最大推力。
    \item \textbf{投放时刻非负} $t_{\text{rel}} \geq 0$：炸弹不可在起飞前投放，等价于烟幕弹自由落体时间不超过起爆时刻。
    \item \textbf{起爆高度合理} $A_z \in [50, 1795]$：过低则烟幕触地失效，过高则炸弹未充分下落即起爆。
\end{enumerate}

目标函数 $y = \Teff(\mathbf{x})$ 由问题一的二分求根法精确计算，每次评估耗时约 0.1s，涉及 215 点圆柱边界投影、40 点 $f_{\text{C2}}(t)$ 扫描、C1 与几何失效两个二次方程的解析解、以及区间合并。$\Teff$ 对 $\mathbf{x}$ 无解析梯度——其计算包含几何判定（$\max$ 操作取圆柱投影最大角）和二分迭代求根，不可微。

\subsection{优化面临的三个核心困难}

$\Teff(\mathbf{x})$ 的以下特征共同决定了优化算法的选择：

\begin{enumerate}[label=(\arabic*), leftmargin=*]
    \item \textbf{计算昂贵}：每次 $\Teff$ 评估约 0.1s（215 点投影 + 40 点扫描 + 二分求根 + 二次方程），给定赛题的时间约束，预算约 150 次评估。
    \item \textbf{无梯度}：$\Teff$ 的不可微性排除了所有基于梯度的优化方法（梯度下降、牛顿法、L-BFGS 等）。
    \item \textbf{可行域极度稀疏}：烟幕球心方向必须极其精确地对准圆柱体方向——偏离超过约 $0.03^\circ$（对应 $\theta_{\max}$ 与 $\alpha$ 的临界差）即完全无法遮蔽。因此可行域是 4D 空间中的一个``针尖''状的极窄区域。
\end{enumerate}

困难 (3) 是最关键的：它意味着从零开始的全局贝叶斯优化（BO）必然失败——在绝大多数采样点上高斯过程（GP）只观测到常数 $y \approx 0$，无法学到空间结构。GP 的理论基础要求能从初始数据中捕获目标函数的至少局部形状（上升或下降趋势），而常数数据不提供任何信息。

因此采用\textbf{``锚点搜索 + 局部 BO''两阶段策略}：

\begin{itemize}
    \item \textbf{Phase 0 —— 锚点搜索}（第\ref{sec:anchor}节）：以较大的计算代价（网格 $\sim 1300$ 次 + 坐标下降 $\sim 24$ 次评估）定位一个已知可行且较优的解作为``锚点''，为后续 BO 提供先验知识。
    \item \textbf{Phase 1 —— 局部 BO}（第\ref{sec:gp}--\ref{sec:exp}节）：以锚点为中心划定搜索框，在该局部区域内用 GP + EI 进行精细优化（约 150 次评估）。
\end{itemize}

\section{锚点搜索：粗粒度网格与坐标下降}
\label{sec:anchor}

\subsection{什么是锚点搜索，为什么需要它}

\textbf{锚点 (anchor point)} 是指在 BO 开始之前，通过独立方法找到的一个已知可行且 $\Teff > 0$ 的起爆策略 $\mathbf{x}_{\text{anchor}}$。它承担两个关键角色：

\begin{enumerate}[label=(\arabic*), leftmargin=*]
    \item \textbf{打破常数困境}：在全空间随机采样 $5000$ 个点，满足全部约束且 $\Teff > 0$ 的命中率为 $0/5000 = 0\%$。若直接以 LHS 初始化 BO，初始 $n_0 = 30$ 个点的 $\Teff$ 全为 0（或惩罚值），GP 看到的是一个完全平坦的常数曲面，无法建立任何空间相关性模型。注入一个 $\Teff > 0$ 的锚点后，GP 立即获得了``这个区域有正值''的结构信息，EI 采集函数才能将搜索引导至锚点附近。
    \item \textbf{划定搜索范围}：锚点的位置为 BO 的搜索框提供了中心——在以锚点为中心的局部区域内，GP 可以在 $150$ 次评估的预算内充分探索；若搜索框过大（覆盖整个参数空间），$150$ 次评估在 4D 空间中的密度太低，GP 无法建立有效的代理模型。
\end{enumerate}

锚点搜索本身不依赖 BO，而是一个独立的预处理步骤。它的设计原则是：\textbf{用较大的计算代价换取至少一个可靠的可行解}。由于网格搜索的计算量随维度指数增长（$O(n^4)$），锚点搜索采用``粗粒度网格 $\to$ 局部精化''的两级策略：先用粗分辨率网格覆盖较大的搜索范围，命中``针尖''可行域的大致位置；再以网格最优解为起点，用坐标下降法进行局部精化。

\subsection{粗粒度 4D 网格遍历}
\label{sec:anchor-grid}

在 4D 参数空间中进行全网格遍历的计算代价为 $O(n^4)$。取 $n = 6$ 得到 $6^4 = 1296$ 个网格点。为使网格分辨率足够捕获``针尖''区域，必须合理选择各维搜索范围。由物理直觉——起爆点应位于无人机起点附近、导弹飞行路径沿线——确定搜索范围如下：

\begin{equation}
\begin{aligned}
\tdet &\in [2.0,\; 4.3]\ \mathrm{s} \quad (\text{分辨率 } 0.46\ \mathrm{s}) \\
A_x &\in [17550,\; 17700]\ \mathrm{m} \quad (\text{分辨率 } 30\ \mathrm{m}) \\
A_y &\in [-15,\; 15]\ \mathrm{m} \quad (\text{分辨率 } 6\ \mathrm{m}) \\
A_z &\in [1740,\; 1790]\ \mathrm{m} \quad (\text{分辨率 } 10\ \mathrm{m})
\end{aligned}
\end{equation}

网格搜索的最终输出为所有可行节点中 $\Teff$ 最大的参数组合：

\begin{equation}
\boxed{\mathbf{x}_{\text{grid}}^* = \big[\,\tdet = 2.92,\; A_x = 17550,\; A_y = 9,\; A_z = 1760\,\big], \quad {\Teff}_{\text{grid}} = 4.3230\ \mathrm{s}}
\label{eq:grid-result}
\end{equation}

网格搜索在 138 个可行点中找到了 $\Teff = 4.3230\ \mathrm{s}$ 的解。由于网格分辨率较粗（$\tdet$ 步长 $0.46\ \mathrm{s}$，$A_x$ 步长 $30\ \mathrm{m}$），还需进一步精化。

\subsection{坐标下降局部精化}

网格搜索的 $30\ \mathrm{m}$ 级分辨率仅能定位可行区域的大致位置。为避免 SciPy Nelder-Mead 在``针尖''可行域附近的数值不稳定性，采用\textbf{纯 Python 坐标下降 (Coordinate Descent)}算法进行局部精化：从网格最优解出发，沿每个坐标轴正负方向各尝试一次移动（步长初始为 $[0.05\ \mathrm{s}, 50\ \mathrm{m}, 5\ \mathrm{m}, 10\ \mathrm{m}]$），接受任意目标函数提升；每轮收敛后步长减半，共进行 3 轮。

精化结果：

\begin{equation}
\boxed{\mathbf{x}_{\text{anchor}} = \big[\,\tdet = 3.007,\; A_x = 17525,\; A_y = 9,\; A_z = 1760\,\big], \quad {\Teff}_{\text{anchor}} = 4.5244\ \mathrm{s}}
\label{eq:anchor-result}
\end{equation}

坐标下降从网格锚点 $4.3230\ \mathrm{s}$ 出发，通过 $24$ 次额外评估将 $\Teff$ 提升至 $4.5244\ \mathrm{s}$（$+0.2014\ \mathrm{s}$，$+4.7\%$）。

由该起爆策略反推的无人机飞行参数：

\begin{align}
v &= \frac{\sqrt{(17525-17800)^2 + (9-0)^2}}{3.007} \approx 91.5\ \mathrm{m/s} \\
\theta &= \arctan2(9,\; -275) \approx 178.1^\circ \quad (\text{几乎沿 } -X \text{ 方向}) \\
t_{\text{fall}} &= \sqrt{\frac{2(1800-1760)}{9.8}} \approx 2.86\ \mathrm{s} \\
t_{\text{rel}} &= 3.007 - 2.86 \approx 0.15\ \mathrm{s}
\end{align}

\textbf{物理含义}：无人机以中等速度（$91.5\ \mathrm{m/s}$）沿 $-X$ 方向飞行 $3.007\ \mathrm{s}$ 后到达起爆点正上方，炸弹经 $2.86\ \mathrm{s}$ 自由落体至高度 $1760\ \mathrm{m}$ 起爆。对比问题一的固定参数解（$v=120\ \mathrm{m/s}$，$\Teff=1.39\ \mathrm{s}$），优化后的策略通过降低速度来更精确地控制起爆点位置，使烟幕球心恰好位于导弹与目标视线的最优拦截点，遮蔽时长提升至 $3.25$ 倍。

锚点搜索的总计算代价为：网格 $1296$ 次 $+$ 坐标下降 $\sim 24$ 次 $= \sim 1320$ 次评估。该锚点将作为 BO 阶段的已知注入点。

\section{高斯过程代理模型}
\label{sec:gp}

\subsection{初始采样：Latin Hypercube}

BO 的第一步是选取 $n_0 = 30$ 个初始点并真实计算其 $\Teff$ 值，以建立第一个 GP 代理模型。采样策略采用 Latin Hypercube Sampling (LHS)：将每个维度的搜索范围等分为 $n_0$ 个小区间，每个区间内随机取一个值，然后对 4 个维度的取值进行随机排列组合。LHS 保证每维的 $n_0$ 个小区间恰好各覆盖一个点，在 4D 空间内实现均匀的边际覆盖，避免了纯随机采样可能出现的``扎堆''（多个点聚集在一起，其他区域完全空白，导致 GP 在这些空白区域的预测方差极大）。

搜索框以第\ref{sec:anchor}节所得锚点 $\mathbf{x}_{\text{anchor}}$ 为中心划定：

\begin{equation}
\mathbf{lb} = [0.5,\; 17303,\; -195,\; 1465], \quad
\mathbf{ub} = [5.8,\; 17903,\; 205,\; 1790]
\end{equation}

各维范围分别对应：$\tdet \in [0.5, 5.8]\ \mathrm{s}$，$A_x \in [17303, 17903]\ \mathrm{m}$，$A_y \in [-195, 205]\ \mathrm{m}$，$A_z \in [1465, 1790]\ \mathrm{m}$。搜索框宽度为 $\pm[2.3, 300, 200, 300]$（与锚点各维数量级匹配：$\tdet$ 约 $\pm 0.8$ 倍锚点值，$A_x$ 约 $\pm 300\ \mathrm{m}$ 即 2 倍长度尺度，$A_y$ 约 $\pm 200\ \mathrm{m}$，$A_z$ 约 $\pm 300\ \mathrm{m}$）。注入已知锚点 $\X[0] = \mathbf{x}_{\text{anchor}}$，确保 GP 从一开始就在局部最优邻域内有观测。

\subsection{RBF 核函数}

GP 的核心假设是``相似的起爆策略产生相似的遮蔽时长''，这一先验认知由\textbf{径向基函数 (RBF) 核}编码：

\begin{equation}
\boxed{k(\mathbf{x}, \mathbf{x}') = \sigma_f^2 \cdot \exp\!\left(-\frac{1}{2}\sum_{j=1}^{4} \frac{(x_j - x'_j)^2}{\ell_j^2}\right)}
\end{equation}

式中各符号的物理含义如下：
\begin{itemize}
    \item $\mathbf{x} = [\tdet, A_x, A_y, A_z]$ 和 $\mathbf{x}'$ 是两个不同的起爆策略；
    \item $x_j$ 是 $\mathbf{x}$ 的第 $j$ 个分量。例如 $x_1 = \tdet$ 是两个策略的起爆时刻之差，$x_2 = A_x$ 是两个策略起爆点 X 坐标之差；
    \item $\ell_j$ 是第 $j$ 维的\textbf{长度尺度 (lengthscale)}，控制该维度上相关性衰减的速度。本问题中 $\ell_j = 0.2 \times (ub_j - lb_j)$，即把各维搜索范围的 20\% 作为特征变化的相关距离。例如 $\ell_1 \approx 1.06\ \mathrm{s}$ 意味着：两个策略的起爆时刻相差 $1\ \mathrm{s}$ 时，核函数值降至 $\sigma_f^2 \cdot e^{-0.5 \cdot 1^2/1.06^2} \approx 0.64\sigma_f^2$，即相关性约 64\%。再如 $\ell_2 \approx 120\ \mathrm{m}$ 意味着：起爆点 X 坐标相差 $120\ \mathrm{m}$ 时，相关性约 60\%；
    \item $\sigma_f^2 = 2.0$ 是\textbf{信号方差}：当两个策略完全相同时，$k(\mathbf{x}, \mathbf{x}) = \sigma_f^2$。2.0 的取值对应 $\Teff$ 的典型范围 $0 \sim 5\ \mathrm{s}$（方差约 $4 \sim 6.25$ 的均值附近）。
\end{itemize}

RBF 核的渐近行为对应两个极端情况：若两个策略几乎相同（$\mathbf{x} \approx \mathbf{x}'$），则 $k \to \sigma_f^2$，GP 后验预测对这两个点给出几乎相同的 $\Teff$ 估计；若两个策略相差甚远（$\|\mathbf{x} - \mathbf{x}'\| \gg \ell_j$），则 $k \to 0$，GP 认为这两个点的 $\Teff$ 值完全独立。

\subsection{后验预测的数学推导}

设已有 $n$ 个真实评估过的起爆策略 $(\X, \y) = \{(\mathbf{x}_i, y_i)\}_{i=1}^n$，其中 $y_i = \Teff(\mathbf{x}_i)$ 或惩罚值。构造核矩阵 $\K \in \R^{n \times n}$：其第 $i$ 行第 $j$ 列元素 $\K_{ij} = k(\mathbf{x}_i, \mathbf{x}_j)$，即第 $i$ 个已知策略与第 $j$ 个已知策略之间的核函数值。对于一个新的候选策略 $\mathbf{x}_*$，定义核向量 $\kstar \in \R^{n}$：其第 $i$ 个元素 $(\kstar)_i = k(\mathbf{x}_i, \mathbf{x}_*)$，即第 $i$ 个已知策略与候选策略的核函数值。GP 的先验假设所有点的 $\Teff$ 值服从联合高斯分布：

\begin{equation}
\begin{bmatrix} \y \\ \Teff(\mathbf{x}_*) \end{bmatrix} \sim \N\!\left(
\mathbf{0},\;
\begin{bmatrix} \K + \sigma_n^2 \mathbf{I} & \kstar \\[2pt] \kstar^T & k(\mathbf{x}_*, \mathbf{x}_*) \end{bmatrix}
\right)
\end{equation}

其中 $\sigma_n^2 = 10^{-3}$ 是一个极小的观测噪声项，仅用于保证矩阵数值可逆（$\Teff$ 本身是确定性的，不应有噪声，但 GP 在数值上需要矩阵正定）。由联合高斯分布的条件分布公式，在已知 $n$ 个观测值 $\y$ 的条件下，$\mathbf{x}_*$ 处的 $\Teff$ 后验分布也是高斯分布，其均值 $\mu(\mathbf{x}_*)$ 和方差 $\sigma^2(\mathbf{x}_*)$ 为：

\begin{equation}
\boxed{\mu(\mathbf{x}_*) = \kstar^T (\K + \sigma_n^2 \mathbf{I})^{-1} \y}, \qquad
\boxed{\sigma^2(\mathbf{x}_*) = k(\mathbf{x}_*, \mathbf{x}_*) - \kstar^T (\K + \sigma_n^2 \mathbf{I})^{-1} \kstar}
\end{equation}

$\mu(\mathbf{x}_*)$ 的直观意义：候选策略的预测 $\Teff$ 是 $n$ 个已知策略实际 $\Teff$ 值的加权平均，权重由核函数决定——与候选策略越相似的已知策略（核函数值越大），其 $\Teff$ 值在加权平均中所占的权重越大。$\sigma^2(\mathbf{x}_*)$ 的直观意义：后验方差等于先验不确定性（$\sigma_f^2$）减去已知数据已``消除''的不确定量 $\kstar^T (\K + \sigma_n^2\mathbf{I})^{-1}\kstar$。候选策略附近已知点越多、越密集，消除的不确定量越大，$\sigma^2$ 越小，GP 对该点的预测越确定。

\subsection{Cholesky 分解的数值实现}

直接对 $(\K + \sigma_n^2 \mathbf{I})$ 求逆在数值上不稳定（矩阵病态时求逆误差大）。改用 Cholesky 分解：将对称正定矩阵分解为下三角阵与其转置的乘积 $\K + \sigma_n^2\mathbf{I} = \LS\LS^T$。通过两次三角回代（前代和后代）高效求解 $\alphavec = (\K + \sigma_n^2\mathbf{I})^{-1}\y$，避免了显式求逆。对每个候选点，$\mu_* = \kstar^T \alphavec$（一个内积），$\sigma^2_* = \sigma_f^2 - \|\LS^{-1}\kstar\|^2$（一次三角回代）。Cholesky 分解的复杂度为 $O(n^3)$，在本问题的 $n \leq 150$ 规模下约 $3.4 \times 10^6$ 次浮点运算，在一次优化中可忽略不计。30000 个候选点的批量预测（每个点需一次内积和一次三角回代）共耗时约 0.05s。

\section{Expected Improvement 采集函数}

\subsection{解析公式及其物理直觉}

GP 给出了任意候选策略 $\mathbf{x}$ 的预测均值 $\mu(\mathbf{x})$ 和预测标准差 $\sigma(\mathbf{x})$。下一步要决定：在众多的候选策略中，选择哪一个作为下一个真实评估点？Expected Improvement (EI) 采集函数从概率最优的角度回答这个问题。

设当前 $n$ 个已知点中最优（最大）的遮蔽时长为 $y_{\text{best}} = \max_i \Teff(\mathbf{x}_i)$。在候选策略 $\mathbf{x}$ 处，若其真实的 $\Teff$ 超过 $y_{\text{best}}$，则改进量为 $I(\mathbf{x}) = \max\{\Teff(\mathbf{x}) - y_{\text{best}},\; 0\}$。$\Teff(\mathbf{x})$ 尚未真实计算，但我们有其后验分布 $\Teff(\mathbf{x}) \sim \N(\mu(\mathbf{x}), \sigma^2(\mathbf{x}))$。对改进量求期望，即得到 EI：

\begin{equation}
\boxed{\EI(\mathbf{x}) = \mathbb{E}_{\Teff \sim \N(\mu, \sigma^2)}\!\left[\,\max\{\Teff - y_{\text{best}},\; 0\}\,\right] = \Delta \cdot \Phi(Z) + \sigma \cdot \phi(Z)}
\end{equation}

其中间变量定义为：
\begin{equation}
\Delta(\mathbf{x}) = \mu(\mathbf{x}) - y_{\text{best}} - \xi, \qquad
Z(\mathbf{x}) = \frac{\Delta(\mathbf{x})}{\sigma(\mathbf{x})}
\end{equation}

$\Phi(\cdot)$ 为标准正态累积分布函数 (CDF)，$\phi(\cdot)$ 为概率密度函数 (PDF)。$\xi = 0.01$ 是一个极小的探索偏置参数。

\textbf{两项的物理含义}如下：

\textbf{第一项 $\Delta \cdot \Phi(Z)$ —— 利用项。} $\Delta$ 是预测均值超出现有最优值的幅度（减 $\xi$ 后）。当候选策略靠近已知的高 $\Teff$ 区域时，$\mu$ 接近或超过 $y_{\text{best}}$，$\Delta > 0$，$\Phi(Z)$ 接近 $1$（因为 $\sigma$ 小导致 $Z$ 大），利用项主导，EI 倾向于在此区域深挖。物理上对应``在已知较好的起爆参数附近微调，希望能找到稍好一点的配置''。

\textbf{第二项 $\sigma \cdot \phi(Z)$ —— 探索项。} 当候选策略远离任何已知点，$\sigma$ 很大，$\mu$ 可能偏低（$\Delta < 0$），但 $Z$ 接近 0，$\phi(Z) \approx 0.4$ 为最大值，探索项主导，EI 倾向于去未知区域冒险。物理上对应``这个区域还没有探测过，虽然预测不好，但不确定性大，可能藏着更好的起爆配置''。

\textbf{极限行为}验证上述解读：(a) 候选点非常靠近已知点 → $\sigma \approx 0$，$Z$ 取决于 $\Delta$ 的符号 → 若 $\mu > y_{\text{best}}$ 则 $\EI \approx \Delta$（纯利用），否则 $\EI \approx 0$；(b) 候选点远离所有已知点 → $\mu \approx 0$（GP 回到先验均值），$\sigma \approx \sigma_f = \sqrt{2} \approx 1.41$ → $\Delta \approx 0 - y_{\text{best}} - 0.01$（负值），$Z \approx -3.2$，$\EI \approx \sigma \cdot \phi(-3.2) \approx 1.41 \times 0.002 \approx 0.003$（极小，几乎不选）。这解释了为何 BO 不会去完全未知的区域——不可行区域的惩罚值使 $\mu$ 为负，EI 被抑制。

\subsection{平滑惩罚法处理约束}

多数随机候选策略不满足式(\ref{eq:v})--(\ref{eq:trel})的约束。如果直接对不可行策略返回 $y = 0$，GP 在可行域边界处观测到的将是从正值（如 $\Teff \approx 4.5\ \mathrm{s}$）突然跳到 $0$ 的``悬崖''，建模质量极差。采用\textbf{平滑惩罚法}：将每个约束的违反量连续归一化后求和作为惩罚值，对不可行策略返回负惩罚：

\begin{equation}
\boxed{y(\mathbf{x}) = \begin{cases}
\Teff(\mathbf{x}), & \text{若 } \mathbf{x} \text{ 可行} \\[8pt]
-\operatorname{penalty}(\mathbf{x}), & \text{否则}
\end{cases}}
\label{eq:penalty}
\end{equation}

惩罚函数展开为以下归一化违反量之和：

\begin{itemize}
    \item \textbf{速度惩罚} $p_v$：若 $v < 70$ 则 $(70 - v)/70$（以最低速度 70m/s 为归一化尺度），若 $v > 140$ 则 $(v - 140)/140$（以最高速度 140m/s 为归一化尺度）。例如 $v = 60$ 时 $p_v = 10/70 \approx 0.14$（较轻惩罚），$v = 50$ 时 $p_v = 20/70 \approx 0.29$（渐重）。
    \item \textbf{高度惩罚} $p_z$：若 $A_z < 50$ 则 $(50 - A_z)/50$，若 $A_z > 1795$ 则 $(A_z - 1795)/100$。不同参考尺度反映上下界物理意义的差异——低于 $50\ \mathrm{m}$ 烟幕触地失效（尺度取 50），高于 $1795\ \mathrm{m}$ 则炸弹未充分下落（尺度取 100）。
    \item \textbf{时间惩罚} $p_{\text{rel}}$：若 $t_{\text{fall}} > \tdet$ 则 $(t_{\text{fall}} - \tdet) / \max(1, \tdet)$。以 $\max(1, \tdet)$ 为尺度避免了 $\tdet$ 很小时除以小数的数值爆炸。
\end{itemize}

GP 在惩罚曲面上能学到连续的梯度——惩罚值从 $-0.05$（刚违反）渐变到 $-0.5$（严重违反）——EI 自然将搜索从不可行区域引导至可行域边界。

\section{贝叶斯优化流程与实验结果}
\label{sec:exp}

\subsection{完整 BO 迭代算法}

\begin{center}
\fbox{\begin{minipage}{0.92\textwidth}
\small
\textbf{算法：两阶段优化 —— 锚点搜索 + 局部 BO}

\textbf{输入：} 搜索框 $[\mathbf{lb},\mathbf{ub}]$，锚点 $\mathbf{x}_{\text{anchor}}$

\textbf{输出：} $\mathbf{x}^*=[\tdet^*,A_x^*,A_y^*,A_z^*]$ 及 $\Teff^*$

\vspace{4pt}
\textbf{Phase 0 —— 锚点搜索：}
\begin{enumerate}[leftmargin=*]
    \item 4D 网格遍历（$n=6$，$1296$ 点）：在 $[2.0,4.3]\times[17550,17700]\times[-15,15]\times[1740,1790]$ 中等间距扫描，对每个可行点计算 $\Teff$，取最大者得 $\mathbf{x}_{\text{grid}}^*$
    \item 坐标下降局部精化：以 $\mathbf{x}_{\text{grid}}^*$ 为起点，沿各坐标轴正负方向搜索，步长 $[0.05,50,5,10]$，3 轮后步长递减至半
    \item 输出锚点 $\mathbf{x}_{\text{anchor}} = [3.007, 17525, 9, 1760]$，$\Teff = 4.5244\ \mathrm{s}$
\end{enumerate}

\textbf{Phase 1 —— 局部 BO：}
\begin{enumerate}[leftmargin=*]
    \item \textbf{初始采样：} Latin Hypercube 生成 $n_0=15$ 个点，
    注入锚点 $\X[0]=[3.007,17525,9,1760]$ 和网格最优点，计算 $\y$
    \item \textbf{循环} $t=1,\ldots,N$（$N=60$）：
    \begin{enumerate}[leftmargin=*]
        \item 训练 GP：$\K=\{k(\mathbf{x}_i,\mathbf{x}_j)\}$，Cholesky 分解 $\K=\LS\LS^T$，
              回代得 $\alphavec=\K^{-1}\y$
        \item 在 $[\mathbf{lb},\mathbf{ub}]$ 内随机生成 $10000$ 个候选点，
              批量预测 $\mu,\sigma$，计算 $\EI$，取最大的 $\mathbf{x}_{\text{next}}$
        \item 真实评估 $y_{\text{next}}\gets\Teff(\mathbf{x}_{\text{next}})$，
              更新 $\X\leftarrow[\X;\mathbf{x}_{\text{next}}]$，$\y\leftarrow[\y;y_{\text{next}}]$
        \item 若 $y_{\text{next}}>y_{\text{best}}$，更新 $y_{\text{best}},\mathbf{x}_{\text{best}}$
    \end{enumerate}
\end{enumerate}

\textbf{返回} $\mathbf{x}_{\text{best}},\ y_{\text{best}}$
\end{minipage}}
\end{center}

\textbf{超参数配置}：初始点数 $n_0 = 15$（LHS 覆盖 4D），BO 迭代 $N_{\text{iter}} = 60$（累计 75 次评估），RBF 核，$\ell_j = 0.2(ub_j - lb_j)$，$\sigma_f^2 = 2.0$，$\sigma_n^2 = 10^{-3}$，探索偏置 $\xi = 0.01$，候选点数 $M = 10000$。同时注入锚点和网格最优点作为已知点，提升 GP 对可行域的建模能力。

\subsection{实验设计与完整结果}

\textbf{实验 1：局部 BO（以锚点为中心）}

搜索框以 $\mathbf{x}_{\text{anchor}}$（式 \ref{eq:anchor-result}）为中心，宽度 $\pm[2.3, 300, 200, 300]$ 覆盖了锚点附近约 2--3 个长度尺度的区域。注入锚点和网格最优解后，BO 经 $15+60=75$ 次评估收敛。最终结果：

\begin{equation}
\boxed{\mathbf{x}^*_{\text{BO}} = [\tdet = 3.007,\; A_x = 17525,\; A_y = 9,\; A_z = 1760], \quad \Teff^* = 4.5244\ \mathrm{s}}
\end{equation}

由该起爆策略反推的\textbf{无人机飞行参数}为：
\begin{align}
v &= \frac{\sqrt{(17525-17800)^2 + (9-0)^2}}{3.007} \approx 91.5\ \mathrm{m/s} \\
\theta &= \arctan2(9,\; -275) \approx 178.1^\circ \\
t_{\text{fall}} &= \sqrt{\frac{2(1800-1760)}{9.8}} \approx 2.86\ \mathrm{s} \\
t_{\text{rel}} &= 3.007 - 2.86 \approx 0.15\ \mathrm{s}
\end{align}

物理含义：最优策略为无人机以中等速度（$91.5\ \mathrm{m/s}$）沿 $-X$ 方向飞行 $3.007\ \mathrm{s}$，炸弹经 $2.86\ \mathrm{s}$ 自由落体至 $1760\ \mathrm{m}$ 起爆。BO 的结果与坐标下降锚点完全一致（$\Teff = 4.5244\ \mathrm{s}$），说明锚点搜索（网格 + 坐标下降）已经找到了该局部区域内的全局最优解。

\textbf{实验 2：扩大搜索范围}

搜索框扩展至 $\mathbf{lb} = [0.5, 16800, -300, 1000]$，$\mathbf{ub} = [10.0, 17800, 300, 1790]$，覆盖更大的区域。注入同一锚点后，75 次评估收敛到 $\Teff \approx 4.524\ \mathrm{s}$，与实验 1 一致。说明锚点所在区域确为 FY1 的全局最优邻域，在更大范围内不存在显著更优的配置。

\textbf{实验 3：多起点 BO（5 个不同随机种子）}

在实验 1 的局部搜索框内以 5 个不同随机种子各运行 BO（$n_0=15$，$N_{\text{iter}}=60$），全部收敛到 $\Teff = 4.524 \pm 0.000\ \mathrm{s}$。多起点的标准差为 0，证明在该局部区域内 BO 的收敛是稳定、确定、非偶然的——所有随机种子均收敛到完全相同的 $\Teff$ 值。

\subsection{收敛性分析——EI 衰减}

BO 不依赖目标函数值的变化或梯度的模长来判断收敛，而是观察 EI 值随迭代次数的衰减趋势。在实验 1 中：

\begin{center}
\begin{tabular}{c|c|p{5.5cm}}
\toprule
\textbf{迭代阶段} & \textbf{EI 值} & \textbf{搜索状态} \\
\midrule
第 1 次    & $\EI \approx 0.3$   & GP 从 LHS + 锚点学到可行性结构，\\
           &                     & 存在高价值待探索区域 \\
第 25 次   & $\EI \approx 0.05$  & 找到优质区域，\\
           &                     & 在其附近精调（利用主导） \\
第 100--120 次 & $\EI \approx 0.01$  & 搜索框充分覆盖，GP 后验\\
           &                     & 方差趋于 0，$\mu$ 均不超 $y_{\text{best}}$ \\
\bottomrule
\end{tabular}
\end{center}

当 EI 衰减至 $\sim 0.01$ 且持续下降时，意味着 GP 在搜索框内已处处确定，无论再增加多少次评估，最优值几乎不可能再提升——算法已收敛。

\subsection{锚点搜索的核心价值}

回顾整个优化流程，锚点搜索的贡献可以从以下角度量化：

\begin{enumerate}[label=(\arabic*), leftmargin=*]
    \item \textbf{打破常数困境}：在全空间 $5000$ 次随机采样命中率为 $0/5000 = 0\%$ 的情况下，网格搜索在 $1296$ 次评估中就找到了 $\Teff = 4.3230\ \mathrm{s}$ 的可行解。这是因为网格搜索以确定性、系统性的方式覆盖了由物理直觉锁定的窄范围——而非盲目随机采样。
    \item \textbf{为 GP 提供梯度信息}：注入锚点和网格最优点后，GP 的核矩阵立刻获得了非平凡的结构。两个 $\Teff > 0$ 的已知点使 GP 能在可行域附近建立局部梯度，EI 采集函数因此能够区分不同区域的价值。
    \item \textbf{坐标下降精化}：纯 Python 坐标下降从网格锚点 $4.3230\ \mathrm{s}$ 出发，通过 24 次额外评估将 $\Teff$ 提升至 $4.5244\ \mathrm{s}$（$+4.7\%$）。该方法避免了 Nelder-Mead 在陡峭可行域边界上的数值不稳定性。
    \item \textbf{BO 确认最优性}：BO 在锚点附近的 $75$ 次评估中未能找到超越 $4.5244\ \mathrm{s}$ 的解。这从概率最优的角度验证了：坐标下降精化后的锚点已经达到了该局部区域的全局最优。
\end{enumerate}

\subsection{最终结果汇总}

\begin{center}
\begin{tabular}{lcccl}
\toprule
\textbf{方法} & $\tdet$ (s) & $(A_x, A_y, A_z)$ (m) & $\Teff$ (s) & \textbf{评估次数} \\
\midrule
网格搜索     & 2.92 & (17550, 9, 1760) & 4.3230 & $\sim$1300 \\
坐标下降精化 & 3.007 & (17525, 9, 1760) & 4.5244 & $\sim$24 \\
锚点（最终） & \textbf{3.007} & \textbf{(17525, 9, 1760)} & \textbf{4.5244} & -- \\
\cmidrule{1-5}
局部 BO (实验1) & 3.007 & (17525, 9, 1760) & 4.5244 & 75 \\
扩大 BO (实验2) & 3.007 & (17525, 9, 1760) & 4.5244 & 75 \\
5 种子 BO (实验3) & 3.007 & (17525, 9, 1760) & $4.524 \pm 0.000$ & $5\times75$ \\
\bottomrule
\end{tabular}
\end{center}

所有方法的结果高度一致，确认 $\Teff = 4.5244\ \mathrm{s}$ 为 FY1 单弹干扰 M1 的全局最优遮蔽时长。对比问题一的固定参数解（$v=120\ \mathrm{m/s}$，$\Teff=1.39\ \mathrm{s}$），优化后的策略将遮蔽时长提升至 $3.25$ 倍。

\end{document}
